% !TeX root = ../report.tex

\section{HƯỚNG DẪN SỬ DỤNG MÃ NGUỒN}
\label{sec:huong-dan-ma-nguon}

\subsection{Đường dẫn và cấu trúc thư mục}
\label{subsec:cau-truc-thu-muc}

Kho mã nguồn GitHub chính thức của dự án: \url{https://github.com/Tricancode/CO3054-IoT-BTL-FireAlarm}.

Cấu trúc thư mục thực tế của kho mã nguồn được tổ chức theo bốn phân hệ chính cùng tệp điều phối khởi chạy tại thư mục gốc như Đoạn mã \ref{lst:folder-tree}.

\begin{lstlisting}[caption={Cấu trúc thư mục mã nguồn thực tế của dự án \texttt{CO3054-IoT-BTL-FireAlarm}}, label={lst:folder-tree}]
CO3054-IoT-BTL-FireAlarm/
|-- run_system.py                         # Trinh khoi chay tong the he thong
|-- HUONG_DAN_CHAY.md                     # Huong dan van hanh kich ban A & B
|-- KICH_BAN_QUAY_DEMO.md                 # Kich ban demo & anh minh chung
|-- gateway_hardware_buffer.db            # CSDL SQLite WAL luu dem ngoai tuyen
|-- 01_Firmware_EndNode/                  # Firmware ESP32-S3 YoloUNO (PlatformIO)
|   |-- platformio.ini                    # Cau hinh board esp32-s3-devkitc-1
|   |-- include/tinyml_fire_model.h       # Model TinyML xuat ra C header (T04)
|   |-- include/lv_conf.h                 # Cau hinh LVGL chuan bi mo rong
|   |-- src/main.cpp                      # Cam bien, TFT, Servo/Buzzer/LED, UART
|   `-- tinyml/train_tinyml.py            # Cong cu huan luyen & xuat model (T05)
|-- 02_Local_Gateway/                     # Local IoT Gateway & Edge AI
|   |-- hardware_gateway/LocalGateway.py  # Cau noi UART <-> MQTT QoS 1 (T08, T15)
|   |-- ui/gateway_infrastructure.py      # SCADA Web 5050 & SQLite WAL (T10, N06)
|   |-- ui/isolation_detector.py          # AI Isolation Forest (T09)
|   |-- dashboard/docker-compose.yml      # Cau hinh ThingsBoard tb-postgres (N07)
|   |-- dashboard/deploy_thingsboard.py   # Script bat/tat Docker ThingsBoard
|   |-- dashboard/setup_thingsboard_dashboard.py # Cap phat thiet bi & Dashboard (T11)
|   |-- dashboard/thingsboard_dashboard.json     # Luoc do 10 widget ThingsBoard 3.x
|   |-- dashboard/verify_telemetry.py     # Kiem thu gui telemetry MQTT
|   |-- run_hardware_gateway.py           # Trinh chay rieng cho Gateway phan cung
|   |-- start_system.py                   # Trinh chay ho tro --mode sim/hardware
|   `-- test_system.py                    # Bo kiem thu don vi tu dong (unittest)
|-- 03_CoreIOT_Config/                    # Cau hinh Rule Chain & do hieu nang
|   |-- rule_chains/fire_alarm_rule_chain.json   # Rule Chain 5 nut chuan (T12)
|   |-- rule_chains/rule_chain_guide.md   # Tai lieu mo ta Rule Chain
|   `-- scripts/
|       |-- setup_rule_chain.py           # Tai tu dong Rule Chain qua REST API
|       |-- test_rulechain_alert.py       # Kiem thu logic Rule Chain (6 test)
|       `-- measure_latency.py            # Do do tre vong lap E2E MQTT -> RPC
|-- 04_Docs_Media/                        # Bao cao LaTeX, slide & so lieu
|   |-- doc/report.tex                    # Tai lieu bao cao chinh
|   `-- experiments/latency_results.csv   # Du lieu do do tre thuc nghiem
`-- captured_picture/                     # 12 anh minh chung (T04..T12, N06, N07)
\end{lstlisting}

\subsection{Yêu cầu môi trường}
\label{subsec:yeu-cau-moi-truong}

Để biên dịch, nạp firmware và vận hành đầy đủ hệ thống, cần chuẩn bị môi trường phần cứng và phần mềm như sau:
\begin{itemize}[leftmargin=*, itemsep=2pt]
  \item \textbf{Phần cứng:} 01 bo mạch OhStem YoloUNO (ESP32-S3), cáp USB Type-C truyền dữ liệu, cảm biến MQ-2, cảm biến lửa Flame, DHT11, màn hình TFT SPI ILI9341 $320 \times 240$, động cơ Servo, còi Buzzer, đèn LED và 01 máy tính cá nhân (Windows 10/11 hoặc Linux/Raspberry Pi OS).
  \item \textbf{Môi trường biên dịch nhúng:} Visual Studio Code kèm tiện ích \textbf{PlatformIO IDE} (hoặc PlatformIO Core CLI) để tự động tải toolchain \texttt{espressif32} và 5 thư viện Arduino khai báo trong \texttt{01\_Firmware\_EndNode/platformio.ini} (\texttt{Adafruit GFX Library}, \texttt{Adafruit ILI9341}, \texttt{ESP32Servo}, \texttt{DHT sensor library}, \texttt{Adafruit Unified Sensor}).
  \item \textbf{Môi trường Python:} Python phiên bản \texttt{3.9} trở lên kèm các thư viện bên thứ ba: \texttt{pyserial}, \texttt{paho-mqtt}, \texttt{requests}, \texttt{scikit-learn}, \texttt{numpy}. (Máy chủ Web cục bộ cổng \texttt{5050} và cơ sở dữ liệu đệm SQLite sử dụng trực tiếp thư viện chuẩn của Python).
  \item \textbf{Container hóa (cho Kịch bản A -- Docker Localhost):} Docker Desktop kèm Docker Compose V2 (hoặc V1) để chạy ảnh \texttt{thingsboard/tb-postgres:latest}.
\end{itemize}

\subsection{Cài đặt và cấu hình}
\label{subsec:cai-dat-cau-hinh}

Thực hiện tuần tự các bước cài đặt trên cửa sổ PowerShell hoặc Terminal tại thư mục làm việc:

\begin{enumerate}[leftmargin=*, itemsep=3pt]
  \item \textbf{Tải mã nguồn và cài đặt thư viện Python:}
  Sao chép kho mã nguồn từ GitHub và cài đặt các gói phụ thuộc bằng \texttt{pip} (\texttt{pyserial paho-mqtt requests scikit-learn numpy}).
  \item \textbf{Khởi chạy máy chủ ThingsBoard cục bộ (Kịch bản A -- Docker Desktop đang mở):}
  Chạy \texttt{python} với tệp \path{02_Local_Gateway/dashboard/deploy_thingsboard.py}. Kịch bản tự động gọi \texttt{docker compose up -d} và chờ dịch vụ Spring Boot cùng PostgreSQL sẵn sàng trên cổng HTTP \texttt{9090} và MQTT \texttt{1883} (thời gian khởi tạo lần đầu khoảng $45$--$60\text{ s}$). Tài khoản mặc định gồm Tenant Administrator (\path{tenant@thingsboard.org} / \texttt{tenant}) và System Administrator (\path{sysadmin@thingsboard.org} / \texttt{sysadmin}).
  \item \textbf{Tự động cấp phát thiết bị Gateway và Dashboard 10 widget:}
  Chạy \texttt{python} với tệp \path{02_Local_Gateway/dashboard/setup_thingsboard_dashboard.py}. Script tự động chờ ThingsBoard sẵn sàng tối đa $180\text{ s}$, đăng ký thiết bị \path{FireAlarm_Gateway} kèm Access Token và tạo Dashboard công nghiệp.
  \item \textbf{Tự động nạp Rule Chain báo cháy:}
  Chạy \texttt{python} với tệp \path{03_CoreIOT_Config/scripts/setup_rule_chain.py} kèm tham số \texttt{--host localhost --port 9090} (hoặc thêm cờ \texttt{--dry-run} để kiểm tra tính hợp lệ tệp JSON trước khi đẩy lên máy chủ).
  \item \textbf{Cấu hình thông số linh hoạt (bảo mật thông tin xác thực):}
  Khi triển khai thực tế trên máy chủ đám mây \texttt{app.coreiot.io}, người dùng truyền trực tiếp địa chỉ broker và Access Token riêng thông qua biến môi trường (\texttt{TB\_REST\_HOST}, \texttt{TB\_DEVICE\_TOKEN}, \texttt{COREIOT\_BROKER}, \texttt{COREIOT\_TOKEN}) hoặc cờ dòng lệnh \texttt{--broker}, \texttt{--token} để tránh lộ mã xác thực sản xuất trong mã nguồn.
\end{enumerate}

\subsection{Biên dịch, nạp và vận hành}
\label{subsec:bien-dich-van-hanh}

Đoạn mã \ref{lst:run-commands} tổng hợp đầy đủ các lệnh từ bước huấn luyện lại mô hình TinyML, biên dịch và nạp firmware lên bo mạch YoloUNO ESP32-S3, vận hành hệ thống theo Kịch bản A (Docker Localhost) hoặc Kịch bản B (Cloud \texttt{app.coreiot.io}), đến thực thi các bộ kiểm thử và đo độ trễ tự động.

\begin{lstlisting}[language=bash, caption={Các lệnh cài đặt, nạp firmware, vận hành hệ thống và kiểm thử tự động}, label={lst:run-commands}]
# 1. Clone repo va cai dat thu vien Python
git clone https://github.com/Tricancode/CO3054-IoT-BTL-FireAlarm.git
cd CO3054-IoT-BTL-FireAlarm
pip install pyserial paho-mqtt requests scikit-learn numpy

# 2. (Tuy chon) Huan luyen lai TinyML va xuat tinyml_fire_model.h (T05)
python 01_Firmware_EndNode/tinyml/train_tinyml.py

# 3. Bien dich va nap firmware cho ESP32-S3 YoloUNO (PlatformIO)
cd 01_Firmware_EndNode
pio run -e yolo_uno -t upload
cd ..

# 4. Kich ban A: Khoi chay ThingsBoard Docker local, cap phat Dashboard & RuleChain
python 02_Local_Gateway/dashboard/deploy_thingsboard.py
python 02_Local_Gateway/dashboard/setup_thingsboard_dashboard.py
python 03_CoreIOT_Config/scripts/setup_rule_chain.py --host localhost --port 9090

# 5. Khoi chay toan bo he thong (tu do cong COM va mo ca 2 Dashboard)
python run_system.py
# Hoac Kich ban B: Ket noi thang toi Cloud CoreIOT (khong can Docker)
python run_system.py --broker app.coreiot.io --token <YOUR_DEVICE_TOKEN>
# Hoac chay Gateway API cho thiet bi con YOLOUNO_FIRE_NODE_01 (T15)
python 02_Local_Gateway/hardware_gateway/LocalGateway.py --broker localhost --gateway-api

# 6. Chay bo kiem thu tu dong va do do tre vong lap E2E
python 02_Local_Gateway/test_system.py
python 03_CoreIOT_Config/scripts/test_rulechain_alert.py
python 03_CoreIOT_Config/scripts/measure_latency.py --trials 10
\end{lstlisting}

\textbf{Thứ tự bật và kiểm tra hoạt động:} (1) Bật container Docker ThingsBoard trước khi khởi chạy Gateway; (2) Cắm cáp USB Type-C vào bo mạch ESP32-S3; (3) Chạy \texttt{python run\_system.py} và quan sát trình duyệt tự động mở hai tab \texttt{http://localhost:5050} (Local SCADA Web Monitor) và \texttt{http://localhost:9090} (ThingsBoard Industrial Dashboard); (4) Xác nhận huy hiệu \texttt{End Node: Online}, \texttt{CoreIOT: Connected} và dữ liệu cảm biến cập nhật liên tục mỗi $1\text{ s}$. Nhấn \texttt{Ctrl+C} tại cửa sổ dòng lệnh để dừng an toàn toàn bộ tiến trình con và giải phóng cổng COM.
